% Copyright 2026 Alan Szmyt. All rights reserved.
\section{System design}
\label{sec:system-design}

Antidote is designed as an inspectable research instrument, not as an opaque
audio recommender or an autonomous affect controller. Its unit of inquiry is a
person encountering a particular acoustic realization in a particular moment
after approving the context and journey that informed it. The architecture
therefore preserves each transformation---source record, working projection,
semantic intent, journey plan, model conditioning, generated artifact,
measured acoustics, exposure, response, and any later update---as a separately
identified object. This section addresses RQ1 by defining that chain and the
authority, timing, uncertainty, and failure boundaries around it.

The equations below are a design language. Each is visibly classified as
\emph{definitional}, \emph{conceptual}, a \emph{proposed operationalization},
an \emph{estimated model}, or a \emph{future empirical model}. A definition
names records without asserting a causal law. A conceptual equation transfers
a useful structure from another field. A proposed operationalization states a
quantity that could be implemented and tested. An estimated model produces an
uncertain system view rather than ground truth. A future empirical model is
neither fitted nor used by the present prototype. Appendix~\ref{app:notation-glossary}
centralizes the notation, and Appendix~\ref{app:equation-classification}
records the status and prohibited inference of every equation.

\subsection{Person--moment conditional model}
\label{sec:conceptual-conditional-model}

Let $i$ index an individual and $t$ a decision point within or between
sessions. Antidote does not represent the person as a fixed feature vector.
Instead, it defines the information available for one bounded moment as the
heterogeneous record
\begin{equation}
  z_{i,t} =
  \left(
    o^{\mathrm{usr}}_{i,t},
    o^{\mathrm{opt}}_{i,t},
    c^{\mathrm{work}}_{i,t},
    g_{i,t},
    k^{\mathrm{usr}}_{i,t},
    \mathcal{H}_{i,<t}
  \right).
  \label{eq:moment-record}
\end{equation}
\noindent\textit{ANT-EQ-001---Definitional; implementation: represented in
the executable mock slice.}

Here $o^{\mathrm{usr}}$ contains explicit check-in observations,
$o^{\mathrm{opt}}$ contains any separately consented optional observations,
$c^{\mathrm{work}}$ is the reviewed context projection, $g$ is the desired
journey, $k^{\mathrm{usr}}$ contains direct semantic-control settings, and
$\mathcal{H}$ denotes eligible prior evidence. Equation~\ref{eq:moment-record}
is a tuple, not a sum: its members need not share units and none is silently
converted into ``the person's state.'' Missing optional observations remain
missing. The same individual may therefore produce different moment records
at two times, and the same audio may be associated with different responses
without either response being treated as anomalous.

The desired journey $g_{i,t}$ is a direction and time-indexed proposal---for
example, staying with, softening, exploring, releasing, focusing, or closing---
rather than a promised terminal state. The current schemas permit compact
valence, arousal, intensity, and confidence fields, but preserve free-text
description because a low-dimensional affect coordinate cannot exhaust the
person's meaning. Likewise, prior history is eligible evidence only when its
retention and learning permissions remain active and the person accepts its
relevance to the present moment.

\subsection{Inspectable record chain}
\label{sec:design-record-chain}

The context boundary separates retained source records from replaceable
interpretations. Following the repository's personal-model contract, the
context state is defined as
\begin{equation}
  S^{\mathrm{ctx}}_{i,t} =
  \left(L_{i,t}, P_{i,t}, V_{i,t}\right),
  \label{eq:context-store}
\end{equation}
\noindent\textit{ANT-EQ-002---Definitional; implementation: append-only events
and derived projections exist for the executable mock slice.}

where $L$ is the consented source-event record, $P$ contains derived
projections, and $V$ contains their versions and provenance. This partition is
compatible with local-first and event-derived views, but neither local storage
nor an append-only history proves privacy, truth, or correct interpretation
\citep{kleppmann2019localfirst,lin2026scroll}. A current consent grant
$A_{i,t}$ then constrains the view used for a declared purpose:
\begin{equation}
  c^{\mathrm{work}}_{i,t}
  = \pi\!\left(S^{\mathrm{ctx}}_{i,t}; A_{i,t}\right).
  \label{eq:working-projection}
\end{equation}
\noindent\textit{ANT-EQ-003---Proposed operationalization; implementation:
manual session-scoped projection is enforced in the executable mock slice.}

The projection operator $\pi$ may select, minimize, summarize, or omit records,
but it must preserve source-event identifiers, derivation version, purpose,
review state, expiry, and any person-authored edit or exclusion. Revoked,
expired, excluded, or out-of-purpose fields do not enter the working view. The
source events remain distinguishable from the projection so a new derivation
can be audited without rewriting what was originally authorized. Consent to
inspect, project, generate, retain, learn, export, and publish remains separate;
manual semantic entry remains usable without historical context.

The record chain continues from the accepted working projection to an immutable
moment record, a versioned journey-plan revision, an approved generation
specification, a generation attempt, a content-hashed artifact and feature
report, an actual exposure, and separated response observations. No skipped
link is inferred from a later one: generating an artifact does not establish
playback, and playback does not establish a response.

\AntidoteFigure{consent-scoped-context-projection}

\subsection{Semantic intent, mixins, and adapter compilation}
\label{sec:design-semantic-intent}

Semantic intent is the model-independent statement of what the person wants a
journey to explore. It combines a desired transition with words, metaphors,
memories, inclusions, exclusions, priorities, and bounded acoustic preferences.
A semantic mixin $m_j$ is a typed and versioned control object with a
human-readable purpose, neutral position, bounds, conflicts, and adapter
mapping. It is not a hidden prompt fragment and does not claim a linear mapping
to acoustics or felt emotion.

For mixin $j$, the effective normalized control is determined by an explicit
authority order:
\begin{equation}
  k^{\mathrm{eff}}_{j,t} =
  \begin{cases}
    k^{\mathrm{usr}}_{j,t}, & \text{when the person sets a current value},\\
    k^{\mathrm{acc}}_{j,t}, & \text{when the person accepts a suggestion},\\
    k^{0}_{j}, & \text{otherwise},
  \end{cases}
  \qquad k^{\mathrm{eff}}_{j,t}\in[0,1].
  \label{eq:effective-knob}
\end{equation}
\noindent\textit{ANT-EQ-004---Proposed operationalization; implementation:
the generalized mixin control surface is not implemented.}

An inferred value has no branch in Equation~\ref{eq:effective-knob}: it gains
authority only through explicit acceptance. A current person-authored value
therefore outranks an accepted suggestion, which outranks the neutral default.
An active exclusion or safety constraint may reject a value or block a
transition, but it may not silently substitute a different intent. Stop,
revocation, and exclusion act immediately and are never smoothed.

For ordinary non-safety controls, a visible rate parameter may reduce jitter:
\begin{equation}
  \alpha_{j,t}
  = (1-\rho_j)\alpha_{j,t-1}+\rho_j k^{\mathrm{eff}}_{j,t},
  \qquad 0<\rho_j\leq 1.
  \label{eq:mixin-smoothing}
\end{equation}
\noindent\textit{ANT-EQ-005---Proposed operationalization; implementation:
control smoothing is not implemented or perceptually validated.}

Both $k$ and $\alpha$ are dimensionless. The rate $\rho_j$, neutral zone, and
maximum rate of change are mixin-specific, visible policy rather than model
invention. Smoothing can reduce control jitter; it cannot prove a smooth
musical transition or a helpful experience.

The resulting semantic condition for journey stage $n$ is
\begin{equation}
  q_{i,n} =
  \operatorname{Resolve}\!\left(
    q^{\mathrm{base}}_{i,n}
    \mathbin{\oplus}
    \bigoplus_{j=1}^{J}\alpha_{j,n}m_j;
    \Xi_{i,t},\prec_{i,t}
  \right),
  \label{eq:semantic-condition}
\end{equation}
\noindent\textit{ANT-EQ-006---Definitional; implementation: the typed
generalized composition remains formalized only.}

where $\oplus$ denotes typed composition, $\Xi$ contains active exclusions and
constraints, and $\prec$ records the visible precedence order used to resolve
conflicts. The operator cannot be interpreted as numerical addition or raw
text concatenation. Human-scale controls may be grounded in established
conditioning dimensions such as tempo, key, chord, and beat, while remaining
separate from the model's native prompt representation
\citep{melechovsky2024mustango,tal2024jasco}.

The accepted journey plan $J_{i,t}$ remains model independent. A replaceable
adapter $a$ compiles one stage against the adapter capability set
$\mathcal{K}_a$:
\begin{equation}
  \begin{aligned}
    \gamma^{(a)}_{i,n}
      &= \operatorname{Compile}_{a}
         \!\left(J_{i,t},q_{i,n},\mathcal{K}_a\right),\\
    \varphi^{(a)}_{i,n}
      &= \operatorname{Prompt}_{a}\!\left(\gamma^{(a)}_{i,n}\right).
  \end{aligned}
  \label{eq:adapter-compilation}
\end{equation}
\noindent\textit{ANT-EQ-007---Proposed operationalization; implementation:
only the immutable mock generation boundary is executable.}

The conditioning state $\gamma$ records model-native controls, conflicts,
downgrades, exclusions, uncertainty, and provenance. The prompt projection
$\varphi$ is merely
one adapter-specific projection of $\gamma$; it is not the journey plan or the
canonical semantic record. An unsupported control must produce a visible
downgrade or failure before generation. The person may accept the revised
specification, change the plan, choose another compatible adapter, or stop.
This separation permits model comparison without rewriting Antidote's domain
language and recognizes that requested conditions are not guaranteed to be
realized by current generators \citep{melechovsky2024mustango,tal2024jasco}.

\AntidoteFigure{semantic-intent-mixer}

\subsection{Two-rate architecture and authority}
\label{sec:design-two-rate-architecture}

The proposed architecture uses two clocks because model deliberation and audio
rendering have different timing and authority requirements. A slower
\emph{deliberative loop} operates at meaningful decision points. It may inspect
authorized observations, maintain alternative state hypotheses, combine
person-authored intent, propose a short journey horizon, and request future
generation. Every material change remains visible and correctable. A faster
\emph{deterministic audio loop} schedules already verified material, applies
bounded gain and transition curves, tracks playback, and responds immediately
to stop or underrun conditions. It does not infer psychology or rewrite
semantic intent. Control-rate and render-rate separation has normative and
neighboring-system precedents, but does not establish Antidote's conformance or
response quality \citep{w3c2024webaudio,hutchings2020adaptive,
ehrlich2019closedloop}.

Authority crosses these loops asymmetrically. The person chooses context,
desired journey, mixins, plan acceptance, generation, playback, retention,
export, and whether evidence may inform a later proposal. The Rust domain core
enforces those transitions and owns the canonical session record. A model
worker receives only an approved generation specification and cannot read the
database, change the personal model, initiate playback, or publish an
interpretation. The renderer may halt or fall back when a deadline or safety
condition fails, but cannot silently choose a new emotional goal.

At the source revision identified in this manuscript's build provenance, an
executable synthetic slice already composes a React/Tauri interface, Rust
session authority, session-scoped context review, deterministic Level-1
journey planning, SQLite events, a supervised mock worker, verified synthetic
WAV production, deliberate playback state, response capture, and safety halts.
This is repository evidence, not human-outcome evidence. The generalized
Semantic Intent Mixer, probabilistic state estimation, receding-horizon
controller, generate-ahead queue, continuity engine, real music model, measured
control adherence, privacy-reviewed export, optional sensors, and autonomous
adaptation are not implemented. The governed final-figure specification
requires both groups to be labeled so the proposed system cannot borrow
credibility from the mock path.

``Closed loop'' therefore names a governed possibility for later evidence to
inform a new proposal. In v0, planning is rule guided, responses are recorded,
and a personal-model change requires a separate evidence-linked proposal and
human acceptance. The worker and controller have no autonomous learning
authority.

\AntidoteFigure{holistic-two-rate-architecture}

\subsection{State belief, receding horizon, and generate-ahead buffer}
\label{sec:design-predictive-horizon}

Real-time evaluation begins with observations rather than a verdict about what
the person needs. Let $o_{i,1:t}$ contain only provenance-linked explicit
reports, control changes, playback actions, and optional signals authorized for
the current purpose. A candidate state estimator may retain alternatives as a
belief distribution
\begin{equation}
  b_{i,t}(x)
  = \Pr\!\left(
      X_{i,t}=x
      \mid o_{i,1:t},u_{i,1:t-1}
    \right).
  \label{eq:belief-state}
\end{equation}
\noindent\textit{ANT-EQ-008---Estimated model; implementation: no state estimator or sensor path exists in v0.}

Equation~\ref{eq:belief-state} transfers belief-state notation from planning
under partial observation; Antidote does not claim to implement a POMDP or to
observe a hidden true emotion \citep{kaelbling1998pomdp}. The state space,
observation model, calibration, and confidence representation remain study
questions. Person correction changes the system record and outranks the model
view. Conflicting, missing, stale, or out-of-distribution observations should
widen uncertainty or suspend a proposal rather than create false precision.
Reviews of passive monitoring and electrodermal affect models reinforce this
validation burden: missingness, feature variation, small samples, limited
generalizability, and weaker valence inference constrain any future sensor use
\citep{deangel2022passive,damelio2025eda}.

At a decision point, a conceptual receding-horizon planner may compare a
sequence $U=(u_t,\ldots,u_{t+H_p})$ of semantic and musical control proposals:
\begin{equation}
  \begin{aligned}
    \mathcal{U}^{\mathrm{adm}}_{i,t}
      &= \mathcal{U}\!\left(A_{i,t},\Xi_{i,t},\mathcal{K}_a\right),\\
    U^{*}_{i,t:t+H_p}
      &\in \operatorname*{arg\,min}_{U\in\mathcal{U}^{\mathrm{adm}}_{i,t}}
      \sum_{\tau=t}^{t+H_p}
      \left[
        \bar{L}_{g}(\hat{x}_{i,\tau\mid t},g_{i,\tau},u_{i,\tau})
        +\lambda_{\Delta}\bar{D}_{\Delta}(u_{i,\tau},u_{i,\tau-1})
        +\lambda_{R}\bar{R}(u_{i,\tau})
        +\lambda_{B}\bar{C}_{\mathrm{burden}}(u_{i,\tau})
      \right].
  \end{aligned}
  \label{eq:receding-horizon}
\end{equation}
\noindent\textit{ANT-EQ-009---Conceptual; implementation: v0 uses a
deterministic rule-guided planner rather than this optimizer.}

The admissible set is constrained by active consent $A$, exclusions and safety
rules $\Xi$, and adapter capabilities $\mathcal{K}_a$. The four barred terms
are normalized, dimensionless design costs: represented journey mismatch,
control change, uncertainty or risk, and interaction burden. Their weights are
therefore dimensionless policy choices, not learned clinical utilities.
$\hat{x}_{i,\tau\mid t}$ is a planning estimate, not a predicted lived state.
Only the first proposed action is shown; nothing is committed until the person
approves it, after which the system observes and may replan. This is an
MPC-inspired transfer under JITAI-like decision points, not a stability,
optimality, intervention, or affect-dynamics claim
\citep{garcia1989mpc,nahumshani2018jitai}.

Generation must also finish before its audio is needed. If $B_{i,t}$ is the
verified playable buffer in seconds, the proposed readiness condition is
\begin{equation}
  B_{i,t} \geq
  \ell_{\mathrm{observe}}+\ell_{\mathrm{plan}}+\ell_{\mathrm{generate}}
  +\ell_{\mathrm{verify}}+\ell_{\mathrm{schedule}}+M_{i,t}.
  \label{eq:buffer-deadline}
\end{equation}
\noindent\textit{ANT-EQ-010---Proposed operationalization; implementation:
generation latency and a generate-ahead buffer are not yet benchmarked.}

Every term in Equation~\ref{eq:buffer-deadline} is measured in seconds;
$M_{i,t}$ is an uncertainty margin derived from the declared hardware,
adapter, and latency distribution. If the condition fails, the renderer uses a
previously approved continuation, loop, ambient bed, or graceful ending. It
does not insert unverified bytes or block the audio callback. Emerging
streaming-generation work makes changing-input, chunk-wise generation a
plausible future adapter, but does not establish local performance or
production readiness \citep{wang2026streaming}.

\AntidoteFigure{predictive-horizon-buffer}

\subsection{Semantic, musical, and waveform continuity}
\label{sec:design-continuity}

A continuously changing prompt does not by itself produce a continuous
listening experience. Antidote therefore separates three questions: whether
neighboring semantic conditions change within their declared bounds, whether
the realized audio is musically and acoustically compatible at a boundary, and
whether the renderer joins verified samples without a signal discontinuity.

For adjacent stages $n$ and $n+1$, semantic-plan distance is
\begin{equation}
  D^{\mathrm{sem}}_n
  = d_{\mathcal{Q}}\!\left(q^{\mathrm{tail}}_{i,n},
                           q^{\mathrm{head}}_{i,n+1}\right),
  \label{eq:semantic-continuity}
\end{equation}
\noindent\textit{ANT-EQ-011---Proposed operationalization; implementation:
no validated semantic-distance function is selected.}

where $d_{\mathcal{Q}}$ is a versioned, dimensionless distance over compatible
typed controls and stage roles. It must expose which mixins contributed to the
distance and cannot assume that interpolation in embedding or prompt space is
perceptually linear.

Measured acoustic-boundary distance is independent:
\begin{equation}
  D^{\mathrm{ac}}_n
  = \sum_{r\in\mathcal{F}}
    w_r\,
    d_r\!\left(
      f_r(a^{\mathrm{tail}}_n),
      f_r(a^{\mathrm{head}}_{n+1})
    \right),
  \qquad \sum_{r\in\mathcal{F}}w_r=1.
  \label{eq:acoustic-continuity}
\end{equation}
\noindent\textit{ANT-EQ-012---Proposed operationalization; implementation:
the mock analyzer does not validate musical boundary continuity.}

The feature family $\mathcal{F}$ may include normalized distances for tempo and
beat phase, key or chord compatibility, loudness, spectral balance and timbre,
density, and phrase-boundary alignment. The $w_r$ are nonnegative,
dimensionless, and versioned; the feature-specific $d_r$ functions normalize
otherwise incompatible units before aggregation. Thresholds for
$D^{\mathrm{sem}}$ and $D^{\mathrm{ac}}$ require perceptual and technical
evaluation and may differ by person or listening context. Loudness standards
can anchor measurement practice without supplying a universal transition or
hearing-safety threshold \citep{ebu2023r128}.

Within an accepted overlap window $s\in[0,1]$, one waveform-level candidate is
the equal-power crossfade
\begin{equation}
  y(s)
  = \cos\!\left(\frac{\pi s}{2}\right)a_n(s)
  + \sin\!\left(\frac{\pi s}{2}\right)a_{n+1}(s).
  \label{eq:equal-power-crossfade}
\end{equation}
\noindent\textit{ANT-EQ-013---Proposed operationalization; implementation:
no adaptive continuity renderer exists in v0.}

Both input signals and $y$ share the same sample-amplitude units. The equation
can reduce a gain discontinuity; it cannot repair incompatible meter, phase,
harmony, phrase, timbre, meaning, or emotional pacing. Resequencing,
reorchestration, continuation generation, stem mixing, and a graceful ending
are separate strategies. A transition is eligible for playback only after its
artifact integrity, buffer margin, semantic condition, acoustic measurements,
and relevant safety constraints pass their own checks. Failure of any one
check invokes an approved fallback rather than being hidden by a crossfade.

\AntidoteFigure{semantic-waveform-continuity}

\subsection{Exposure, response, and advisory updating}
\label{sec:design-response-update}

The generated segment $a_n$ is not the exposure $e_{i,t}$. Exposure records
which content hash was actually played, when playback began, duration heard,
completion, interruption, stop reason, and safety events. Response is then a
heterogeneous record linked to that exposure:
\begin{equation}
  \mathsf{Y}_{i,t} =
  \left(
    y^{\mathrm{perc}}_{i,t},
    y^{\mathrm{felt}}_{i,t},
    y^{\mathrm{int}}_{i,t},
    w^{\mathrm{int}}_{i,t},
    y^{\mathrm{help}}_{i,t},
    y^{\mathrm{harm}}_{i,t},
    y^{\mathrm{res}}_{i,t},
    y^{\mathrm{mis}}_{i,t},
    y^{\mathrm{after}}_{i,t}
  \right).
  \label{eq:response-record}
\end{equation}
\noindent\textit{ANT-EQ-014---Definitional; implementation: the mock slice
captures a bounded subset and keeps exposure separate.}

The components describe perceived musical expression, felt state, intensity,
whether that intensity was wanted, immediate helpfulness, possible harm,
resonance, mismatch, and a later aftereffect. They may contain text, bounded
scores, categorical values, or missingness, so Equation~\ref{eq:response-record}
is a record declaration rather than a numeric vector. Intense crying, for
example, cannot be classified as success or harm from intensity alone; the
person's appraisal, stop behavior, and later account remain separate evidence.

A future within-person response model could estimate
\begin{equation}
  p_{\theta_i}\!\left(
    \mathsf{Y}_{i,t}
    \mid e_{i,t},z_{i,t}
  \right),
  \label{eq:response-model}
\end{equation}
\noindent\textit{ANT-EQ-015---Future empirical model; implementation: no
response predictor is fitted or used for generation.}

where $\theta_i$ parameterizes a person-specific and versioned hypothesis.
Conditioning on exposure prevents an unplayed artifact from becoming response
evidence. The model would need to represent uncertainty, changed meanings,
missing responses, contradictory sessions, and alternative explanations; it
could not infer clinical benefit from acoustic adherence or immediate affect.

If a future protocol justified Bayesian updating, a candidate posterior form
would be
\begin{equation}
  p\!\left(\theta_i\mid\mathcal{D}_{i,1:t}\right)
  \propto
  p\!\left(
    \mathsf{Y}_{i,t}
    \mid e_{i,t},z_{i,t},\theta_i
  \right)
  p\!\left(\theta_i\mid\mathcal{D}_{i,<t}\right).
  \label{eq:personal-update}
\end{equation}
\noindent\textit{ANT-EQ-016---Future empirical model; implementation: v0
records evidence-linked proposals and prohibits autonomous updating.}

Equation~\ref{eq:personal-update} is deliberately generic: it does not specify
a valid likelihood, prior, forgetting rule, causal effect, or sufficient sample
size. In the current authority model, eligible observations may support a
descriptive proposal, but a new personal-model snapshot is created only after
the person reviews the evidence and accepts the update. Rejection preserves
both the prior snapshot and the rejected proposal. The design follows an
interactive-machine-learning rationale while retaining the possibility that
feedback requests can be burdensome or reduce trust
\citep{amershi2014interactive,honeycutt2020feedback}.

\AntidoteFigure{moment-journey-equation-map}

\subsection{Provenance and failure boundaries}
\label{sec:design-provenance-failures}

Every derived record links to the activity, agent, version, and source records
that produced it. For generation this includes the accepted plan hash, adapter
and model identity, code revision, effective parameters, seed when supported,
device class, elapsed time, warnings, downgrade decisions, output hashes, and
feature-analyzer versions. For exposure and response it includes the artifact
hash, playback interval, interruption state, observation window, instrument
version, missingness, corrections, and permission for later use. W3C PROV
supplies an interoperable vocabulary for entities, activities, agents, and
derivations, while Workflow Run RO-Crate demonstrates a packaging pattern for
computational lineage \citep{moreau2013prov,leo2024rocrate}. Antidote proposes
to use those concepts; a provenance record remains an assertion to audit, not
proof that its interpretation is true.

The system fails closed at the boundary whose precondition is missing. Invalid,
expired, or ambiguous consent prevents projection. Conflicting mixins or an
unsupported capability prevent compilation until the conflict or downgrade is
reviewed. A worker crash, malformed response, timeout, partial output, path
violation, size violation, or hash mismatch cannot become a playable verified
artifact. A missed generation deadline invokes an approved fallback. A stop,
reported harm, or safety event halts automatic continuation and is never scored
as engagement. A failed or rejected model update preserves the prior snapshot.
An export remains blocked when classification, minimization, redaction,
consent, privacy review, or destination approval is incomplete.

Append-only events make those transitions reconstructable, while separately
classified payloads permit future retention and deletion policies to act on
sensitive content. The present repository implements integrity and replay for
the synthetic path; encryption, key recovery, backup, secure deletion,
projection invalidation, real-model sandboxing, and privacy-reviewed research
export remain unresolved. Local execution and content hashes reduce some
dependencies and improve traceability, but they do not establish security,
ownership, meaningful consent, safety, or benefit.

\AntidoteFigure{provenance-privacy-export}
